\makeabstract
    {
    Sixty Years of Ethnic Studies: A Scoping Review of Emerging Research on Ethnic Studies Policy and Impacts on Adolescent Development 
    }
    {
    Kevin Romero\textsuperscript{1}, Angel Fu\textsuperscript{1},  Ava Stanton\textsuperscript{1}, Miguel Kamgaing\textsuperscript{2}, Sarah Gillespie\textsuperscript{1}, Andres Pinedo\textsuperscript{2}
    }
    {
    \textsuperscript{1} Department of Psychology, Amherst College, Amherst, MA\\
\textsuperscript{2} Peabody College of Education and Human Development, Vanderbilt University, Nashville, TN
    }
    {
    Amidst response to student protests for an education that reflects the experiences and perspectives of communities of color, ethnic studies was created to examine our nation{\textquoteright}s diverse racial and ethnic groups through social, economic, and historical lenses (Sleeter \& Zavala, 2020, p. iv). But in recent years, amidst social justice demonstrations and divided legislative proposals to expand or ban ethnic studies, it is important, now more than ever, to understand the implications of ethnic studies for adolescent development. Our scoping review will address the following questions: How are schools defining and implementing ethnic studies? What mechanisms have been identified as the {\textquotedblleft}active ingredients{\textquotedblright} of ethnic studies? And what has research shown about the implications of ethnic studies for youth{\textquoteright}s development? With advances in both theoretical and empirical research, we aim to provide an updated review of ethnic studies in U.S. middle and high schools to inform policy and practice.
    }
    {
    Our review followed PRISMA guidelines for scoping reviews (Tricco et al., 2018) to synthesize the most recent trends in ethnic studies implementation and research. Results were imported into Covidence software for automatic deduplication, screening, and data extraction. All titles and abstracts (n=4236) and full texts (n=132) were screened by two of the four research assistants using a criteria list, achieving high group agreement (\ensuremath{>}95\%). ProQuest was used to extract works from two databases (PsycINFO and ERIC) from January 2019 to March 2026.
    }
    {
    Ethnic Studies classes can bolster success across a range of developmental tasks in a diverse society. As we continue data extraction, with planned completion in Fall 2026, there are preliminary results to highlight. More articles on ethnic studies have been published in the past six years than in the previous 40 years (1980s{\textendash}2020) covered in Sleeter \& Zavala{\textquoteright}s (2020) scoping review. And just as the field of ethnic studies continues to expand, so do the research methods being used to understand it, ranging from large-scale quantitative analysis to autoethnographic research by ethnic studies teachers themselves. And as the U.S. K{\textendash}12 student body continues to diversify, teachers are reporting innovative pedagogies they use to integrate ethnic studies in the classroom.
    }
    {
    Almost 60 years ago, student protestors at San Francisco State College sparked a new era for U.S. education, one focused on implementing the experiences, voices, and perspectives of those who have been marginalized in U.S. history. As we synthesize the results of our research, we plan to submit our scoping review to the Special Issue in the Journal of Research on Adolescence in February 2027. And with our results, we hope to translate our findings for others through knowledge translation, including school officials, teachers, policymakers, parents, community members, and, most importantly, those who helped make what we now know as the field of Ethnic Studies: students.
    }

    \newpage
    